\chapter*{Заключение}\label{ux437ux430ux43aux43bux44eux447ux435ux43dux438ux435}
\addcontentsline{toc}{chapter}{Заключение}

\begin{spacing}{1.2}
\fontsize{13}{16.5}\selectfont
В курсовой работе рассмотрены методы машинного обучения для генерации названий растений и продуктов, а также разработано программное средство, объединяющее генеративную модель и модули оценки результатов. Теоретический анализ показал, что для небольшой специализированной предметной области рационально использовать компактную LSTM-модель: она учитывает контекст последовательности, обучается на центральном процессоре и не требует значительных вычислительных ресурсов.

Для обучения сформирован корпус названий растений и продуктов, проведены очистка, нормализация и удаление повторов. Дополнительно применена аугментация за счёт сочетания названий с прилагательными, что позволило увеличить объём обучающих данных примерно втрое. Для посимвольной LSTM реализовано условие по категории и длине названия; использованы два слоя, скрытое состояние размером 128, dropout 0,2 и оптимизатор Adam. Время обучения на CPU составляет около одной-двух минут.

Для отбора результатов разработаны три модуля скоринга: классификатор стиля на основе char-CNN, оценщик правдоподобия на основе биграммной языковой модели и оценщик новизны на основе n-граммного сравнения. В совокупности они ранжируют кандидатов и помогают выбирать варианты, соответствующие заданной категории. В ходе тестирования перплексия LSTM составила около 4--5 на валидационной выборке, а точность классификатора стилей достигла 80\%. Сравнение с марковской цепью и n-граммной моделью показало, что LSTM формирует более плавные и осмысленные названия.

Разработанный интерфейс на Gradio содержит отдельные вкладки для растений и продуктов. Пользователь может задать категорию, длину названия и количество вариантов, сохранить понравившиеся результаты в избранное и экспортировать их в текстовый или CSV-файл. Таким образом, поставленная цель достигнута: создано локальное приложение, пригодное для генерации названий в аграрных симуляторах, играх и прикладных задачах нейминга. Дальнейшую работу можно направить на расширение датасета, развитие словесных моделей с учётом морфологии и использование пользовательской обратной связи.
\end{spacing}
